\documentclass[12pt,letterpaper]{article}
\usepackage[utf8]{inputenc}
\usepackage[spanish,es-noshorthands]{babel}
\usepackage{amsmath,amssymb}
\usepackage{graphicx}
\usepackage{xcolor}
\usepackage{tikz}
\usetikzlibrary{calc,arrows.meta,decorations.markings}
\usepackage[top=2.5cm, bottom=2.7cm, left=2cm, right=2cm, footskip=40pt]{geometry}
\usepackage{fancyhdr}
\usepackage{lastpage}

\definecolor{ubbblue}{RGB}{0,51,102}

\pagestyle{fancy}
\fancyhf{}
\renewcommand{\headrulewidth}{0pt}
\renewcommand{\footrulewidth}{0.4pt}
\fancyfoot[L]{\raisebox{-0.5\height}{\includegraphics[height=1.1cm]{logo ubb.png}}}
\fancyfoot[C]{%
    \begin{minipage}[c]{8cm}\centering
        \footnotesize\textcolor{ubbblue}{\textbf{Dr.\ Arturo Fern\'andez P\'erez}}\\
        \scriptsize Departamento de F\'isica -- Universidad del B\'io-B\'io\\
        \scriptsize \textit{Electro\_Asistente\_AI-UBB}\\
        \tiny P\'agina \thepage\ de \pageref{LastPage}
    \end{minipage}}
\fancyfoot[R]{\raisebox{-0.5\height}{\includegraphics[height=1.1cm]{logo dfis.png}}}
\fancypagestyle{plain}{%
    \fancyhf{}
    \renewcommand{\headrulewidth}{0pt}
    \renewcommand{\footrulewidth}{0.4pt}
    \fancyfoot[L]{\raisebox{-0.5\height}{\includegraphics[height=1.1cm]{logo ubb.png}}}
    \fancyfoot[C]{%
        \begin{minipage}[c]{8cm}\centering
            \footnotesize\textcolor{ubbblue}{\textbf{Dr.\ Arturo Fern\'andez P\'erez}}\\
            \scriptsize Departamento de F\'isica -- Universidad del B\'io-B\'io\\
            \scriptsize \textit{Electro\_Asistente\_AI-UBB}\\
            \tiny P\'agina \thepage\ de \pageref{LastPage}
        \end{minipage}}
    \fancyfoot[R]{\raisebox{-0.5\height}{\includegraphics[height=1.1cm]{logo dfis.png}}}
}

\begin{document}

%========================================
% ENCABEZADO INSTITUCIONAL
%========================================
\begin{center}
    \includegraphics[height=2cm]{logo ubb.png}\\[0.3cm]
    {\large\textcolor{ubbblue}{\textbf{UNIVERSIDAD DEL B\'IO-B\'IO}}}\\[0.1cm]
    {\normalsize Facultad de Ciencias -- Departamento de F\'isica}
\end{center}

\vspace{0.3cm}
\hrule height 1pt
\vspace{0.5cm}

{\centering\Large\textbf{Problemas Resueltos en la Clase del 21/10}\par}
{\centering\normalsize\textit{Fuerza Magn\'etica sobre Cargas y Corrientes}\par}

\vspace{0.5cm}

%--- recordatorio: productos cruz de unitarios ---
\begin{center}
\fcolorbox{ubbblue}{ubbblue!5}{%
\begin{minipage}{0.92\textwidth}
\centering
\textbf{\textcolor{ubbblue}{Productos cruz de los vectores unitarios cartesianos}}\\[4pt]
$\hat{\imath}\times\hat{\jmath}=\hat{k}$,\quad
$\hat{\jmath}\times\hat{k}=\hat{\imath}$,\quad
$\hat{k}\times\hat{\imath}=\hat{\jmath}$
\qquad\qquad
$\hat{\jmath}\times\hat{\imath}=-\hat{k}$,\quad
$\hat{k}\times\hat{\jmath}=-\hat{\imath}$,\quad
$\hat{\imath}\times\hat{k}=-\hat{\jmath}$\\[3pt]
$\hat{\imath}\times\hat{\imath}=\hat{\jmath}\times\hat{\jmath}=\hat{k}\times\hat{k}=\vec{0}$
\end{minipage}}
\end{center}

\vspace{0.4cm}

%========================================
% PROBLEMA 1: FUERZA SOBRE UN PROTÓN
%========================================
\section*{Problema 1: Fuerza magn\'etica sobre un prot\'on}

\noindent Un haz de protones se mueve con una rapidez $v = 3\times 10^{5}\ [\text{m/s}]$ a trav\'es
de un campo magn\'etico uniforme de magnitud $B = 2\ [\text{T}]$. Seg\'un la figura, el campo
$\vec{B}$ es \textbf{paralelo al eje $z$} y la velocidad $\vec{v}$ forma un \'angulo de
$30^\circ$ con respecto al eje $z$, en el plano $x$--$z$. Encuentre la fuerza magn\'etica sobre
un prot\'on.

\medskip

\noindent
\begin{minipage}[t]{0.40\textwidth}
\centering
\vspace{0pt}
\begin{tikzpicture}[scale=1.0, >=Stealth, line width=0.9pt]
    % plano x-z (horizontal) sombreado
    \fill[gray!12] (-0.1,0.27)--(2.69,0.27)--(0.09,-1.53)--(-2.71,-1.53)--cycle;
    % ejes
    \draw[->,thick] (0,0)--(2.5,0) node[right]{$x$};
    \draw[->,thick] (0,0)--(0,2.2) node[above]{$y$};
    \draw[->,thick] (0,0)--(-1.82,-1.26) node[below left]{$z$};
    % proton
    \shade[ball color=red!70] (0,0) circle (5pt);
    \node[above right=0pt,font=\small] at (0.06,0.12){$q$};
    % B (azul) paralelo al eje z
    \draw[->,blue,line width=1.6pt] (0,0)--(-1.40,-0.97)
          node[above left=-1pt,font=\small]{$\vec{B}$};
    % v (verde) a 30° del eje z, en el plano x-z (inclinada hacia +x)
    \draw[->,green!60!black,line width=1.6pt] (0,0)--(-0.67,-1.56)
          node[below=1pt,font=\small]{$\vec{v}$};
    % angulo 30° entre el eje z (B) y v
    \draw[gray,line width=0.6pt] (214.7:0.72) arc(214.7:246.8:0.72);
    \node[font=\footnotesize] at (-0.66,-0.66){$30^\circ$};
    % F resultante: -y (perpendicular al plano x-z, hacia abajo)
    \draw[->,orange!85!black,line width=1.4pt] (0,0)--(0,-1.65);
    \node[right=1pt,font=\small,text=orange!85!black] at (0.03,-1.35){$\vec{F}$};
\end{tikzpicture}
\end{minipage}%
\hfill
\begin{minipage}[t]{0.56\textwidth}
\vspace{0pt}
\textbf{Datos:}
\begin{itemize}\setlength\itemsep{1pt}
    \item $q = +1{,}6\times10^{-19}\ [\text{C}]$ (prot\'on)
    \item $v = 3\times10^{5}\ [\text{m/s}]$,\quad $B = 2\ [\text{T}]$
    \item \'Angulo entre $\vec{v}$ y el eje $z$: $\;\theta = 30^\circ$
\end{itemize}

\smallskip
\textbf{Forma vectorial de la fuerza magn\'etica:}
\[
    \vec{F} = q\,\vec{v}\times\vec{B}
\]

\smallskip
Seg\'un el sistema de referencia, $\vec{B}$ es paralelo al eje $z$ y
$\vec{v}$ forma $30^\circ$ con el eje $z$ en el plano $x$--$z$:
\begin{align*}
    \vec{v} &= v\sin30^\circ\,\hat{\imath} + v\cos30^\circ\,\hat{k} \\
    \vec{B} &= B\,\hat{k}
\end{align*}
\end{minipage}

\bigskip

\noindent Reemplazando en la forma vectorial y distribuyendo el producto cruz:
\begin{align*}
    \vec{F} &= q\,\bigl(v\sin30^\circ\,\hat{\imath} + v\cos30^\circ\,\hat{k}\bigr)\times\bigl(B\,\hat{k}\bigr) \\[4pt]
            &= q\,v\,B\sin30^\circ\,(\hat{\imath}\times\hat{k})
             + q\,v\,B\cos30^\circ\,(\hat{k}\times\hat{k})
\end{align*}

\noindent Evaluando \textbf{cada producto cruz de unitarios por separado}:
\[
    \hat{\imath}\times\hat{k} = -\hat{\jmath},
    \qquad\qquad
    \hat{k}\times\hat{k} = \vec{0}
\]

\noindent Por lo tanto:
\begin{align*}
    \vec{F} &= -\,q\,v\,B\sin30^\circ\,\hat{\jmath} + \vec{0} \\[4pt]
            &= -\,(1{,}6\times10^{-19})(3\times10^{5})(2)\,(0{,}5)\,\hat{\jmath}
\end{align*}

\[
    \boxed{\vec{F} = -\,4{,}8\times10^{-14}\,\hat{\jmath}\ [\text{N}]}
\]

\noindent\small\textit{Interpretaci\'on:} la fuerza es perpendicular al plano que contiene a
$\vec{v}$ y $\vec{B}$ (el plano $x$--$z$), es decir, a lo largo del eje $y$, en sentido
$-\hat{\jmath}$. Su magnitud es $F = qvB\sin30^\circ = 4{,}8\times10^{-14}\ [\text{N}]$.
\normalsize

\newpage

%========================================
% PROBLEMA 2: LEVITACIÓN MAGNÉTICA DE UN CABLE
%========================================
\section*{Problema 2: Levitaci\'on magn\'etica de un cable conductor}

\noindent Un cable transporta una corriente $I = 50\ [\text{A}]$ de \textbf{Oeste a Este}, en una
regi\'on con un campo magn\'etico uniforme de magnitud $B = 2\ [\text{T}]$ dirigido
$45^\circ$ \textbf{al Norte del Este} (en el plano horizontal). Determine la fuerza magn\'etica
sobre una secci\'on de $L = 1\ [\text{m}]$ de cable y la masa que debe tener el cable de cobre
para levitar (su peso equilibrado por la fuerza magn\'etica).

\medskip

\noindent
\begin{minipage}[t]{0.42\textwidth}
\centering
\vspace{0pt}
\begin{tikzpicture}[scale=1.0, >=Stealth, line width=0.9pt]
    % rosa de los vientos (vista superior del plano horizontal)
    \draw[->,gray!70] (0,0)--(0,2.0) node[above,black,font=\small]{Norte};
    \draw[->,gray!70] (0,0)--(0,-1.9) node[below,black,font=\small]{Sur};
    \draw[->,gray!70] (0,0)--(2.5,0) node[right,black,font=\small]{Este};
    \draw[->,gray!70] (0,0)--(-2.5,0) node[left,black,font=\small]{Oeste};
    % cable Oeste-Este
    \draw[orange!75!black,line width=3pt] (-2.2,0)--(2.2,0);
    \draw[->,black,line width=1.4pt] (-0.7,0)--(0.9,0)
          node[midway,below=2pt,font=\small]{$I$};
    % B a 45° N del E
    \draw[->,blue,line width=1.7pt] (0,0)--(1.45,1.45)
          node[above right=0pt,font=\small]{$\vec{B}$};
    \draw[gray,line width=0.6pt] (0.6,0) arc(0:45:0.6);
    \node[font=\footnotesize] at (0.92,0.30){$45^\circ$};
    % F sale del plano (+k)
    \draw[black,line width=0.9pt] (1.5,-1.0) circle(7pt);
    \fill[black] (1.5,-1.0) circle(1.4pt);
    \node[below,font=\footnotesize,align=center] at (1.5,-1.25){$\vec{F}$ sale\\del plano};
\end{tikzpicture}\\[3pt]
{\footnotesize $\hat{\imath}$: Este,\; $\hat{\jmath}$: Norte,\; $\hat{k}$: vertical (hacia arriba).}
\end{minipage}%
\hfill
\begin{minipage}[t]{0.54\textwidth}
\vspace{0pt}
\textbf{Datos:}
\begin{itemize}\setlength\itemsep{1pt}
    \item $I = 50\ [\text{A}]$,\quad $L = 1\ [\text{m}]$,\quad $B = 2\ [\text{T}]$
    \item Corriente: Oeste $\to$ Este $\;\Rightarrow\; \vec{L} = L\,\hat{\imath}$
    \item $\vec{B}$ a $45^\circ$ al Norte del Este (plano $x$--$y$)
\end{itemize}

\smallskip
\textbf{Forma vectorial de la fuerza sobre un conductor recto:}
\[
    \vec{F} = I\,\vec{L}\times\vec{B}
\]

\smallskip
Vectores en el plano horizontal:
\begin{align*}
    \vec{L} &= L\,\hat{\imath} = (1)\,\hat{\imath} \\
    \vec{B} &= B\cos45^\circ\,\hat{\imath} + B\sin45^\circ\,\hat{\jmath}
\end{align*}
\end{minipage}

\bigskip

\noindent Reemplazando en la forma vectorial y distribuyendo el producto cruz:
\begin{align*}
    \vec{F} &= I\,(L\,\hat{\imath})\times\bigl(B\cos45^\circ\,\hat{\imath} + B\sin45^\circ\,\hat{\jmath}\bigr) \\[4pt]
            &= I\,L\,B\cos45^\circ\,(\hat{\imath}\times\hat{\imath})
             + I\,L\,B\sin45^\circ\,(\hat{\imath}\times\hat{\jmath})
\end{align*}

\noindent Evaluando \textbf{cada producto cruz de unitarios por separado}:
\[
    \hat{\imath}\times\hat{\imath} = \vec{0},
    \qquad\qquad
    \hat{\imath}\times\hat{\jmath} = \hat{k}
\]

\noindent Por lo tanto:
\begin{align*}
    \vec{F} &= \vec{0} + I\,L\,B\sin45^\circ\,\hat{k} \\[4pt]
            &= (50)(1)(2)\left(\tfrac{\sqrt{2}}{2}\right)\hat{k}
\end{align*}

\[
    \boxed{\vec{F} = 70{,}71\,\hat{k}\ [\text{N}]}
\]

\noindent La fuerza apunta verticalmente hacia arriba ($+\hat{k}$, fuera del plano horizontal).

\newpage

%========================================
% LEVITACIÓN MAGNÉTICA (página aparte)
%========================================
\section*{Condici\'on de levitaci\'on magn\'etica}

\noindent Inicialmente el cable descansa sobre el piso. Para que el cable \textbf{se despegue del
suelo}, la fuerza normal $N$ que el piso ejerce sobre \'el debe \textbf{anularse} ($N = 0$).
Analicemos las fuerzas que act\'uan sobre el cable a lo largo del eje vertical $z$:

\medskip

\noindent
\begin{minipage}[t]{0.34\textwidth}
\centering
\vspace{0pt}
\begin{tikzpicture}[scale=1.0, >=Stealth, line width=0.9pt]
    % eje z
    \draw[->,gray!55] (0,-2.0)--(0,2.4) node[above,black,font=\small]{$z$};
    % cable (corte transversal)
    \shade[ball color=orange!70] (0,0) circle(6pt);
    \node[right=3pt,font=\footnotesize] at (0.18,0){cable};
    % F magnética (arriba, sobre el cable)
    \draw[->,blue,line width=1.6pt] (0,0)--(0,1.9)
          node[right,font=\small]{$\vec{F}$};
    % N normal (arriba, desplazada a la izquierda)
    \draw[->,gray!45!black,line width=1.4pt] (-0.7,0)--(-0.7,1.2)
          node[left,font=\small]{$\vec{N}$};
    % peso (abajo)
    \draw[->,black,line width=1.6pt] (0,0)--(0,-1.7)
          node[right,font=\small]{$m\vec{g}$};
\end{tikzpicture}\\[3pt]
{\footnotesize Diagrama de cuerpo libre del cable.}
\end{minipage}%
\hfill
\begin{minipage}[t]{0.62\textwidth}
\vspace{0pt}
\textbf{Fuerzas sobre el eje $z$:}
\begin{itemize}\setlength\itemsep{2pt}
    \item Peso: $\;-mg\,\hat{k}$ \hfill (hacia abajo)
    \item Normal del piso: $\;+N\,\hat{k}$ \hfill (hacia arriba)
    \item Fuerza magn\'etica: $\;+F\,\hat{k}$, \; con $F = 70{,}71\ [\text{N}]$ \hfill (hacia arriba)
\end{itemize}

\smallskip
\textbf{Segunda ley de Newton} (equilibrio vertical, $\sum F_z = 0$):
\[
    N + F - mg = 0
    \qquad\Longrightarrow\qquad
    N = mg - F
\]
\end{minipage}

\bigskip

\noindent En el instante en que el cable \textbf{se despega del piso}, la fuerza normal se anula.
Con $N = 0$, las \'unicas fuerzas sobre el eje $z$ son el \textbf{peso} y la \textbf{fuerza
magn\'etica}, que quedan equilibradas:
\[
    N = 0
    \qquad\Longrightarrow\qquad
    F - mg = 0
    \qquad\Longrightarrow\qquad
    F = mg
\]

\noindent Despejando la masa del cable:
\[
    m = \frac{F}{g} = \frac{70{,}71}{9{,}8}
\]

\[
    \boxed{m = 7{,}21\ [\text{kg}]}
\]

\noindent\small\textit{Interpretaci\'on:} esta es la masa m\'axima que puede tener el cable para
que la fuerza magn\'etica logre levantarlo. Si la masa fuera mayor, se tendr\'ia $mg > F$ y el
cable permanecer\'ia apoyado en el piso (con $N > 0$).
\normalsize

\newpage

%========================================
% PROBLEMA 3: LAZO TRIANGULAR (Certamen 3 2021-1, Forma B, P1)
%========================================
\section*{Problema 3: Fuerza magn\'etica sobre un lazo de corriente triangular}

\noindent Un lazo de corriente $I = 1{,}87\ [\text{A}]$ circula en una regi\'on de campo magn\'etico
uniforme de magnitud $B = 500\ [\text{mT}]$, el cual apunta \textbf{hacia afuera de la p\'agina},
como muestra la figura. Considere el sistema de referencia de la figura y que la longitud del tramo
$ac$ es $l_{ac} = 18\ [\text{cm}]$. Determine la fuerza magn\'etica $\vec{F}$ sobre cada uno de los
tres segmentos que forman el lazo de corriente.

\medskip

\noindent
\begin{minipage}[t]{0.44\textwidth}
\centering
\vspace{0pt}
\begin{tikzpicture}[scale=1.0, >=Stealth, line width=0.9pt]
    % campo saliente (puntos rojos)
    \foreach \x in {0.55,1.75,2.95,4.05} \foreach \y in {-0.45,0.75,1.85,3.05}
        \fill[red!75!black] (\x,\y) circle (1.7pt);
    % ejes
    \draw[->,thick] (0,0)--(4.6,0) node[right]{$x$};
    \draw[->,thick] (0,0)--(0,3.7) node[above]{$y$};
    % lazo triangular
    \draw[line width=1.4pt] (0,0)--(3.6,0)--(0,3.0)--cycle;
    \node[below left]  at (0,0)   {$a$};
    \node[below]       at (3.6,-0.05) {$b$};
    \node[left]        at (0,3.0) {$c$};
    % angulo de 40 grados en b
    \draw[gray,line width=0.6pt] ($(3.6,0)+(140:0.7)$) arc(140:180:0.7);
    \node[font=\footnotesize] at ($(3.6,0)+(160:1.05)$) {$40^\circ$};
    % sentido de la corriente a -> b -> c -> a
    \draw[->,black,line width=1.1pt] (0.7,0.28)--(1.7,0.28);
    \draw[->,black,line width=1.1pt] (2.51,0.52)--(1.51,1.35);
    \draw[->,black,line width=1.1pt] (-0.35,2.1)--(-0.35,1.0) node[midway,left]{$I$};
    % fuerzas (naranja), proporcionales a su magnitud
    \draw[->,orange!85!black,line width=1.6pt] (1.8,0)--(1.8,-0.8)
        node[right,font=\small]{$\vec{F}_{ab}$};
    \draw[->,orange!85!black,line width=1.6pt] (1.8,1.5)--({1.8+0.67},{1.5+0.80})
        node[right,font=\small]{$\vec{F}_{bc}$};
    \draw[->,orange!85!black,line width=1.6pt] (0,0.55)--(-0.67,0.55)
        node[left,font=\small]{$\vec{F}_{ca}$};
\end{tikzpicture}\\[3pt]
{\footnotesize Campo saliente (puntos); corriente $a\to b\to c\to a$.}
\end{minipage}%
\hfill
\begin{minipage}[t]{0.52\textwidth}
\vspace{0pt}
\textbf{Datos:}
\begin{itemize}\setlength\itemsep{1pt}
    \item $I = 1{,}87\ [\text{A}]$,\quad $l_{ac} = 0{,}18\ [\text{m}]$
    \item \'Angulo en el v\'ertice $b$: $40^\circ$
    \item Campo saliente: $\vec{B} = 0{,}500\,\hat{k}\ [\text{T}]$
\end{itemize}

\smallskip
\textbf{Longitudes de los tramos} (tri\'angulo rect\'angulo en $a$):
\begin{align*}
    \tan40^\circ = \frac{l_{ac}}{l_{ab}} &\;\Rightarrow\; l_{ab} = \frac{0{,}18}{\tan40^\circ} = 0{,}215\ [\text{m}]\\[2pt]
    \sin40^\circ = \frac{l_{ac}}{l_{bc}} &\;\Rightarrow\; l_{bc} = \frac{0{,}18}{\sin40^\circ} = 0{,}280\ [\text{m}]
\end{align*}

\smallskip
\textbf{Forma vectorial de la fuerza sobre un tramo recto:}
\[
    \vec{F} = I\,\vec{l}\times\vec{B}
\]
\end{minipage}

\bigskip

\noindent\textbf{Vector de cada tramo}, en el sentido de la corriente ($a\to b\to c\to a$):
\begin{align*}
    \vec{l}_{ab} &= 0{,}215\,\hat{\imath}\ [\text{m}]\\[2pt]
    \vec{l}_{bc} &= -\,l_{bc}\cos40^\circ\,\hat{\imath} + l_{bc}\sin40^\circ\,\hat{\jmath}
                  = -0{,}215\,\hat{\imath} + 0{,}18\,\hat{\jmath}\ [\text{m}]\\[2pt]
    \vec{l}_{ca} &= -0{,}18\,\hat{\jmath}\ [\text{m}]
\end{align*}

\noindent\textbf{Tramo $ab$} (de $a$ hacia $b$):
\begin{align*}
    \vec{F}_{ab} &= I\,\vec{l}_{ab}\times\vec{B} = (1{,}87)(0{,}215\,\hat{\imath})\times(0{,}500\,\hat{k})\\[2pt]
                 &= (1{,}87)(0{,}215)(0{,}500)\,(\hat{\imath}\times\hat{k}) = 0{,}2010\,(-\hat{\jmath})
\end{align*}
\[
    \boxed{\vec{F}_{ab} = -\,0{,}2010\,\hat{\jmath}\ [\text{N}]}
\]

\noindent\textbf{Tramo $bc$} (de $b$ hacia $c$), distribuyendo el producto cruz:
\begin{align*}
    \vec{F}_{bc} &= I\,\vec{l}_{bc}\times\vec{B}
                  = (1{,}87)\bigl(-0{,}215\,\hat{\imath} + 0{,}18\,\hat{\jmath}\bigr)\times(0{,}500\,\hat{k})\\[2pt]
                 &= (1{,}87)(0{,}500)\bigl[-0{,}215\,(\hat{\imath}\times\hat{k}) + 0{,}18\,(\hat{\jmath}\times\hat{k})\bigr]\\[2pt]
                 &= 0{,}935\bigl[-0{,}215\,(-\hat{\jmath}) + 0{,}18\,\hat{\imath}\bigr]
\end{align*}
\[
    \boxed{\vec{F}_{bc} = 0{,}1683\,\hat{\imath} + 0{,}2010\,\hat{\jmath}\ [\text{N}]}
\]

\noindent\textbf{Tramo $ca$} (de $c$ hacia $a$):
\begin{align*}
    \vec{F}_{ca} &= I\,\vec{l}_{ca}\times\vec{B} = (1{,}87)(-0{,}18\,\hat{\jmath})\times(0{,}500\,\hat{k})\\[2pt]
                 &= -(1{,}87)(0{,}18)(0{,}500)\,(\hat{\jmath}\times\hat{k}) = -\,0{,}1683\,\hat{\imath}
\end{align*}
\[
    \boxed{\vec{F}_{ca} = -\,0{,}1683\,\hat{\imath}\ [\text{N}]}
\]

\noindent\textbf{Verificaci\'on (fuerza neta).} En un campo uniforme la fuerza neta sobre un lazo
cerrado es nula:
\[
    \vec{F}_{ab} + \vec{F}_{bc} + \vec{F}_{ca}
    = (0{,}1683 - 0{,}1683)\,\hat{\imath} + (-0{,}2010 + 0{,}2010)\,\hat{\jmath} = \vec{0}\ \checkmark
\]

\noindent\small\textit{Interpretaci\'on:} cada fuerza es perpendicular a su tramo y apunta hacia
\emph{afuera} del tri\'angulo. La corriente circula en sentido antihorario, de modo que el momento
magn\'etico del lazo apunta en $+\hat{k}$, paralelo a $\vec{B}$: las fuerzas tienden a estirar el
lazo, pero no lo trasladan ni lo hacen girar. La magnitud de $\vec{F}_{bc}$ es
$I\,l_{bc}\,B = (1{,}87)(0{,}280)(0{,}500) \approx 0{,}26\ [\text{N}]$.
\normalsize

\end{document}
